% Einheit 1 von 3 – Ganzrationale Funktionen (bank/grenzwerte-und-verhalten-im-unendlichen/e1.jsonl, zusammenbau v0.1)
%% TODO Zeitmarke Einheit 1: Marken-Zeile nicht lesbar
%% TODO Zweigzeile Teil 1: was hier gelernt wird (Satz in Schülersprache aus dem Einheitstitel) – Einheit 1
\einheitenkopf[e1]{Einheit 1 von 3 · Ganzrationale Funktionen}
\zweigzeile{Abitur GK}
\setcounter{aufgabe}{9}

%% TODO Titel: Ich-kann-Satz fehlt in der Bank (Platzhalter: Kettenname) – Nr. 10
%% TODO Anweisung: ein Satz über der ersten Teilaufgabe fehlt in der Bank – Nr. 10
\begin{aufgabe}{„Wer gewinnt?“}
\begin{teile}
\teil Wer gewinnt für sehr große $x$? $f(x) = 5x^2 - 0{,}1x^3 + 7$ \kreuz{$5x^2$} \\ \kreuz{$-0{,}1x^3$} \\ \kreuz{$7$}
\end{teile}
\end{aufgabe}

%% TODO Titel: Ich-kann-Satz fehlt in der Bank (Platzhalter: Kettenname) – Nr. 11
%% TODO Anweisung: ein Satz über der ersten Teilaufgabe fehlt in der Bank – Nr. 11
\begin{aufgabe}{Ganzrationale Funktionen}
%% TODO \rechenplatz{n}: Schrittzahl des Grundfalls fehlt in der Bank (nur bei fester Schreibform, 2.3 b)
\begin{teile}
\teil Gerade oder ungerade, plus oder minus? $f(x) = -3x^5 + x^2$ \kreuz{Grad gerade, Koeffizient positiv} \\ \kreuz{Grad gerade, Koeffizient negativ} \\ \kreuz{Grad ungerade, Koeffizient positiv} \\ \kreuz{Grad ungerade, Koeffizient negativ}
\teil $f(x) = x^2 - 6x + 1$. Verhalten für $x \to +\infty$ und $x \to -\infty$? x → +∞: f(x) → \leerfeld ; x → −∞: f(x) → \leerfeld
\teil $f(x) = -0{,}5x^4 + 3x^2 + 2$. Gib das Verhalten der Funktionswerte im Unendlichen an. \hfill (FHR 2023) \\ x → +∞: f(x) → \leerfeld ; x → −∞: f(x) → \leerfeld
\teil $f(x) = 3x^2 - \frac{1}{5}x^5 + 4$. Gib das Verhalten der Funktionswerte im Unendlichen an. \hfill (FHR 2022) \\ x → +∞: f(x) → \leerfeld ; x → −∞: f(x) → \leerfeld
\teil $f(x) = \frac{1}{20}x^5 - x^2$ hat die Ableitung $f'(x) = \frac{1}{4}x^4 - 2x$. Gib das Verhalten von $f'$ für $x \to -\infty$ und $x \to +\infty$ an. \hfill (Abitur 2026 GK) \\ x → +∞: f'(x) → \leerfeld ; x → −∞: f'(x) → \leerfeld
\teil $f(x) = 0{,}5x^4 - 3x^2 + 12$. Begründe, dass der Graph achsensymmetrisch zur $y$-Achse ist. \hfill (FHR 2021)
\teil $f(x) = -\frac{1}{4}x^4 + 3x^2 + 5$; der Punkt $Q(2 | 13)$ liegt auf dem Graphen. Begründe, dass der Graph achsensymmetrisch zur $y$-Achse ist, und gib das Verhalten der Funktionswerte im Unendlichen an. \hfill (FHR 2023)
\end{teile}
\end{aufgabe}

%% TODO Titel: Ich-kann-Satz fehlt in der Bank (Platzhalter: Kettenname) – Nr. 12
%% TODO Anweisung: ein Satz über der ersten Teilaufgabe fehlt in der Bank – Nr. 12
\begin{aufgabe}{Ganzrationale Funktionen – Fehler finden, Begründen, Darstellungswechsel}
\begin{teile}
\teil $f(x) = 2x^5 - x^2$. Ein Schüler schreibt: „Für $x \to +\infty$ und für $x \to -\infty$ gilt $f(x) \to +\infty$.“ Finde den Fehler.
\teil Begründe, warum für sehr große $x$ bei $f(x) = x^4 - 50x^3$ der Summand $x^4$ entscheidet.
\teil Welcher Term passt zum Graphen? \kreuz{$f(x) = -0{,}25x^4 + 2x^2$} \\ \kreuz{$f(x) = 0{,}25x^4 - 2x^2$} \\ \kreuz{$f(x) = -0{,}25x^3 + 2x$}

\begin{ksys}[xmin=-4,xmax=4,ymin=-3,ymax=5,ablesen] \funktion{-0.25*\x^4+2*\x^2}{f} \end{ksys}
\end{teile}
\end{aufgabe}
